\documentclass[letterpaper]{article}
\usepackage[scaled=0.92]{helvet}
\renewcommand{\familydefault}{\sfdefault}
\usepackage{amsmath,amssymb}
\usepackage{geometry}
\usepackage{multicol}
\usepackage{enumitem}
\usepackage{microtype}
\usepackage{titlesec}
\usepackage{mathtools}
\usepackage{array}
\usepackage{booktabs}
\pagestyle{empty}

\geometry{landscape,margin=0.08in}
\DeclareMathSizes{5}{5}{5}{5}
\setlength{\parindent}{0pt}
\setlength{\parskip}{0pt}
\setlength{\columnsep}{0.06in}
\setlength{\abovedisplayskip}{0.6pt}
\setlength{\belowdisplayskip}{0.6pt}
\setlength{\abovedisplayshortskip}{0.4pt}
\setlength{\belowdisplayshortskip}{0.4pt}
\setlist[itemize]{leftmargin=0.6em,itemsep=0pt,topsep=0pt,parsep=0pt,partopsep=0pt}
\titlespacing*{\section}{0pt}{0.8pt}{0.4pt}
\titlespacing*{\subsection}{0pt}{0.6pt}{0.3pt}
\titleformat{\section}{\bfseries\fontsize{5}{5.2}\selectfont}{}{0pt}{}
\titleformat{\subsection}{\bfseries\fontsize{5}{5.2}\selectfont}{}{0pt}{}
\newcommand{\bb}[1]{\mathbb{#1}}
\newcommand{\R}{\bb{R}}
\newcommand{\C}{\bb{C}}
\newcommand{\Span}{\operatorname{span}}
\newcommand{\Col}{\operatorname{Col}}
\newcommand{\Nul}{\operatorname{Nul}}
\newcommand{\Rank}{\operatorname{rank}}
\newcommand{\diag}{\operatorname{diag}}
\newcommand{\norm}[1]{\left\lVert #1\right\rVert}
\newcommand{\ket}[1]{|#1\rangle}
\newcommand{\bra}[1]{\langle #1|}
\newcommand{\braket}[2]{\langle #1 \mid #2\rangle}
\newcommand{\sig}{\sigma}
\newcommand{\lam}{\lambda}
\newcommand{\smallhead}[1]{{\bfseries #1}}

\begin{document}
\fontsize{5}{5.2}\selectfont

\begin{center}
{\bfseries Practice Exam 2.1 Cheat Sheet}\\[-0.2ex]
{Letter paper, landscape, 2 pages for double-sided printing, uniform 5pt text/math sizing.}
\end{center}
\vspace{-0.8ex}

\begin{multicols*}{4}
\section{Spectral Theorem Core}
\smallhead{Symmetric $\Leftrightarrow$ orthogonally diagonalizable.} For real matrices:
\[
A \text{ orthogonally diagonalizable } \iff A^T=A.
\]
If $A=PDP^T$ with $P$ orthogonal, then $A^T=P D^T P^T=A$ since $D$ is diagonal.

\smallhead{Spectral Theorem.} If $A=A^T$, then $A$ has an orthonormal eigenbasis, all eigenvalues are real, and
\[
A=PDP^T,\quad P^TP=I,\quad D=\diag(\lambda_1,\dots,\lambda_n).
\]
Columns of $P$ are orthonormal eigenvectors.

\smallhead{Orthogonality of eigenspaces.} If $A=A^T$, $Au=\lam u$, $Av=\mu v$, and $\lam\neq \mu$, then
\[
u\cdot v=0.
\]
Proof trick:
\[
\lam(u\cdot v)=(Au)\cdot v=u\cdot(Av)=\mu(u\cdot v).
\]
Hence $(\lam-\mu)(u\cdot v)=0$.

\smallhead{Multiplicity fact.} For symmetric $A$, geometric multiplicity = algebraic multiplicity for each eigenvalue. So if $\lam$ appears $k$ times on diagonal of $D$, then $\dim E_\lam=k$.

\smallhead{Quick proof identities.}
\[
(Ax)\cdot y=(Ax)^Ty=x^TA^Ty=x^T(Ay)=x\cdot(Ay).
\]
If $A$ symmetric and $B$ any matrix:
\[
(B^TAB)^T=B^TAB,\quad (B^TB)^T=B^TB,\quad (BB^T)^T=BB^T.
\]
If $A=PDP^T$ invertible, then
\[
A^{-1}=PD^{-1}P^T.
\]
If $A,B$ symmetric and $AB=BA$, then
\[
(AB)^T=B^TA^T=BA=AB,
\]
so $AB$ is symmetric and therefore orthogonally diagonalizable.

\smallhead{Upper triangular trap.} If $A=PRP^{-1}$ with $P$ orthogonal and $A=A^T$, then
\[
PRP^T=PR^TP^T \Rightarrow R=R^T.
\]
Upper triangular + symmetric $\Rightarrow$ diagonal.

\section{Spectral T/F Traps}
\smallhead{25.} Orthogonally diagonalizable $\Rightarrow$ symmetric: \textbf{true}.\\
\smallhead{26.} Some symmetric not orthogonally diagonalizable: \textbf{false}.\\
\smallhead{27.} Every orthogonal matrix orthogonally diagonalizable: \textbf{false}; need symmetry, not just orthogonality.\\
\smallhead{28.} $PDP^T$ symmetric: \textbf{true}.\\
\smallhead{29.} $vv^T$ projection: \textbf{false unless $v$ unit}. Correct projector onto $\Span\{v\}$ is
\[
\frac{1}{v^Tv}vv^T.
\]
If $u$ is unit, then
\[
(uu^T)^2=u(u^Tu)u^T=uu^T.
\]
\smallhead{30.} Distinct-eigenvalue eigenvectors of symmetric matrix orthogonal: \textbf{true}.\\
\smallhead{31.} Symmetric matrix has $n$ distinct eigenvalues: \textbf{false}.\\
\smallhead{32.} Eigenspace dimension can be smaller than multiplicity: \textbf{false for symmetric matrices}.

\section{SVD Core}
\smallhead{Definition.}
\[
A=U\Sigma V^T
\]
with $U$ orthogonal $m\times m$, $V$ orthogonal $n\times n$, and
\[
\Sigma=\diag(\sigma_1,\dots,\sigma_r,0,\dots),\quad \sigma_1\ge \sigma_2\ge \cdots\ge 0.
\]

\smallhead{How to compute by hand.}
\begin{itemize}
\item Compute $A^TA$.
\item Find eigenvalues $\lambda_i$ of $A^TA$.
\item Singular values are $\sigma_i=\sqrt{\lambda_i}$.
\item Find orthonormal eigenvectors $v_i$ of $A^TA$.
\item For $\sigma_i>0$, compute
\[
u_i=\frac{1}{\sigma_i}Av_i.
\]
\item Complete $U$ by adding orthonormal vectors if needed.
\end{itemize}

\smallhead{Facts from SVD.}
\[
\Rank(A)=\#\{\sigma_i>0\}.
\]
\[
\Col(A)=\Span\{u_1,\dots,u_r\}.
\]
\[
\Nul(A)=\Span\{v_{r+1},\dots,v_n\}.
\]
Largest stretch:
\[
\max_{\norm{x}=1}\norm{Ax}=\sigma_1,
\]
attained at $x=v_1$ (or any first right singular vector).

\smallhead{Square invertible case.}
\[
A^{-1}=V\Sigma^{-1}U^T.
\]
\smallhead{Determinant fact.} If $A$ is square,
\[
|\det A|=\sigma_1\cdots \sigma_n.
\]
Reason: $|\det U|=|\det V|=1$.

\section{2x2 SVD Template}
For
\[
A=\begin{pmatrix}a&b\\c&d\end{pmatrix},
\quad
A^TA=\begin{pmatrix}a^2+c^2&ab+cd\\ab+cd&b^2+d^2\end{pmatrix}.
\]
Then solve
\[
\det(A^TA-\lambda I)=0
\]
to get $\lambda_1,\lambda_2$, hence $\sigma_1,\sigma_2$. Find unit eigenvectors $v_1,v_2$ of $A^TA$. Then compute
\[
u_i=\frac{1}{\sigma_i}Av_i.
\]
Always check:
\[
U^TU=I,\quad V^TV=I,\quad U\Sigma V^T=A.
\]

\section{Image Compression Core}
For image matrix $A$,
\[
A=U\Sigma V^T,\qquad
A_k=\sum_{i=1}^k \sigma_i u_i v_i^T.
\]
$A_k$ is the best rank-$k$ approximation in 2-norm and Frobenius norm.

\smallhead{Energy / quality preserved.}
\[
\text{quality}(k)=\frac{\sigma_1^2+\cdots+\sigma_k^2}{\sigma_1^2+\cdots+\sigma_n^2}.
\]
On the practice exam, this is what you report for $75\%,50\%,25\%$ kept or for fixed $k$.

\smallhead{Storage count.} Original $m\times n$ image stores $mn$ numbers. Rank-$k$ form stores
\[
km+kn+k=k(m+n+1).
\]
So exact compression ratio:
\[
\frac{k(m+n+1)}{mn}.
\]
Approximate ratio:
\[
\frac{k(m+n)}{mn}.
\]
Worth compressing when
\[
k(m+n+1)<mn
\quad\Leftrightarrow\quad
k<\frac{mn}{m+n+1}.
\]
For square $n\times n$ image:
\[
\text{approx ratio}\approx \frac{2k}{n},\quad
\text{worth it if }k<\frac n2.
\]

\smallhead{Exam 10.1 numbers.}
\[
\sigma \approx (2.4503,1.1180,0.8350,0.5464).
\]
Quality:
\[
k=3:96.38\%,\quad k=2:87.93\%,\quad k=1:72.78\%.
\]

\columnbreak
\section{Quantum Bases}
\smallhead{Standard basis.}
\[
\ket{0}=\binom10,\qquad \ket{1}=\binom01.
\]
\smallhead{Hadamard basis.}
\[
\ket{+}=\frac{1}{\sqrt2}(\ket0+\ket1),\qquad
\ket{-}=\frac{1}{\sqrt2}(\ket0-\ket1).
\]
\smallhead{$i$-basis.}
\[
\ket{i}=\frac{1}{\sqrt2}(\ket0+i\ket1),\qquad
\ket{-i}=\frac{1}{\sqrt2}(\ket0-i\ket1).
\]

\smallhead{Same state criterion.} States represent the same physical state iff they differ by a global phase:
\[
\ket{\psi}\sim e^{i\theta}\ket{\psi}.
\]
Examples:
\[
-\ket0\sim \ket0,\quad i\ket1\sim \ket1,\quad i\ket1\sim \ket1.
\]
But relative phase matters:
\[
\frac{1}{\sqrt2}(\ket0+\ket1)\not\sim \frac{1}{\sqrt2}(-\ket0+i\ket1).
\]

\smallhead{Born rule.} If $\{\ket{b_1},\ket{b_2}\}$ is measurement basis and state is $\ket{\psi}$, then
\[
P(\ket{b_j})=|\braket{b_j}{\psi}|^2.
\]
Probabilities always sum to $1$.

\smallhead{Useful expansions.}
\[
\ket0=\frac{1}{\sqrt2}(\ket+ + \ket-),\qquad
\ket1=\frac{1}{\sqrt2}(\ket+ - \ket-).
\]
\[
\frac{1}{\sqrt2}(\ket+ + \ket-)=\ket0,\quad
\frac{1}{\sqrt2}(\ket+ - \ket-)=\ket1.
\]
\[
\frac{1}{\sqrt2}(\ket{i}+\ket{-i})=\ket0,\quad
\frac{1}{\sqrt2}(\ket{i}-\ket{-i})=i\ket1.
\]

\smallhead{Superposition test.} In a given basis, the state is a superposition iff both basis coefficients are nonzero. Example:
\[
\ket+ = \frac{1}{\sqrt2}\ket0+\frac{1}{\sqrt2}\ket1
\]
is a superposition in the standard basis but not in the Hadamard basis.

\smallhead{Fast probability facts.}
\[
P_{\{\ket0,\ket1\}}(\ket0 \text{ from } \tfrac{\sqrt3}{2}\ket0-\tfrac12\ket1)=\tfrac34.
\]
\[
P_{\{\ket+,\ket-\}}(\ket+ \text{ from } \ket0)=\tfrac12.
\]
\[
P_{\{\ket{i},\ket{-i}\}}(\ket{i}\text{ from }\ket1)=\tfrac12.
\]

\section{Quantum Practice Traps}
\smallhead{Choose the right basis to separate states.}
\begin{itemize}
\item If states already look like $\ket+$ vs $\ket-$, measure in $\{\ket+,\ket-\}$.
\item If one state is $\ket{i}$ and another is orthogonal to it, measure in $\{\ket{i},\ket{-i}\}$.
\item If expressions differ only by global phase, they are the same state, so no measurement distinguishes them.
\end{itemize}

\smallhead{Common exam equalities.}
\[
\frac{1}{\sqrt2}(\ket1-\ket0)=-\frac{1}{\sqrt2}(\ket0-\ket1)\sim \ket-.
\]
\[
\frac{1}{\sqrt2}(\ket{i}-\ket{-i})=i\ket1\sim\ket1.
\]
\[
\frac{1}{\sqrt2}(\ket{i}+\ket{-i})=\ket0.
\]
\[
\frac{1}{\sqrt2}(\ket-+\ket+)=\ket0.
\]
\[
\frac{1}{\sqrt2}(\ket0+e^{i\pi/4}\ket1)\sim
\frac{1}{\sqrt2}(e^{-i\pi/4}\ket0+\ket1).
\]

\section{What To Memorize}
\begin{itemize}
\item Symmetric $\Leftrightarrow$ orthogonally diagonalizable.
\item Distinct eigenvalues of symmetric matrix give orthogonal eigenvectors.
\item $A^TA$ gives singular values and right singular vectors.
\item $u_i=Av_i/\sigma_i$.
\item Rank/Col/Nul from SVD: first nonzero-$\sigma$ columns of $U$, last zero-$\sigma$ columns of $V$.
\item Compression ratio $=k(m+n+1)/(mn)$.
\item Same quantum state iff global phase only.
\item Probability = squared magnitude of inner product with basis vector.
\end{itemize}

\end{multicols*}

\newpage
\fontsize{5}{5.2}\selectfont
\begin{center}
{\bfseries Practice Exam 2.1 Cheat Sheet -- Tricky Problems + One Usual Problem}
\end{center}
\vspace{-0.8ex}

\begin{multicols*}{4}
\section{Nonstandard 1: Projection Trap}
\smallhead{Question type.} ``Is $vv^T$ a projection matrix?''\\
\smallhead{Correct move.} Test idempotence:
\[
(vv^T)^2=v(v^Tv)v^T=(v^Tv)vv^T.
\]
This equals $vv^T$ iff $v^Tv=1$. So:
\begin{itemize}
\item nonzero arbitrary $v$: usually \textbf{not} a projection matrix;
\item unit $u$: $uu^T$ \textbf{is} the orthogonal projector onto $\Span\{u\}$.
\end{itemize}
Fast counterexample:
\[
v=\binom21,\quad vv^T=\begin{pmatrix}4&2\\2&1\end{pmatrix},
\]
\[
(vv^T)^2=5vv^T\neq vv^T.
\]

\section{Nonstandard 2: Read Everything From Given SVD}
If exam gives
\[
A=U\Sigma V^T
\]
already, \textbf{do not recompute}. Just read:
\[
\Rank(A)=\#\{\sigma_i>0\}.
\]
\[
\Col(A)=\Span\{u_1,\dots,u_r\}
\]
where $u_i$ are first $r$ columns of $U$.
\[
\Nul(A)=\Span\{v_{r+1},\dots,v_n\}
\]
where these are columns of $V$ corresponding to zero singular values.
If only $V^T$ is shown, remember: \textbf{rows of $V^T$ = transposes of columns of $V$}. So the last row of $V^T$ gives the last column of $V$.

\section{Nonstandard 3: Max-Length Question}
Question type: ``Find a unit vector $x$ where $\norm{Ax}$ is maximal.''\\
Answer immediately:
\[
x=v_1,\qquad \norm{Ax}_{\max}=\sigma_1.
\]
No need to optimize with calculus. Find the first right singular vector.
For Practice Exam 2.1 Exercise 7:
\[
A=\begin{pmatrix}2&-1\\2&2\end{pmatrix},
\quad
v_1=\frac{1}{\sqrt5}\binom21,\quad
\sigma_1=3.
\]
Therefore
\[
\norm{Av_1}=3.
\]

\section{Nonstandard 4: Compression Ratio Derivation}
Question type: ``Using $k$ singular values for an $m\times n$ image, what is the compression ratio? When is it worth doing?''

Store:
\[
u_1,\dots,u_k \in \R^m \Rightarrow km \text{ numbers},
\]
\[
v_1,\dots,v_k \in \R^n \Rightarrow kn \text{ numbers},
\]
\[
\sigma_1,\dots,\sigma_k \Rightarrow k \text{ numbers}.
\]
Total stored:
\[
k(m+n+1).
\]
Original storage:
\[
mn.
\]
Hence
\[
\text{ratio}=\frac{k(m+n+1)}{mn}.
\]
Worth it:
\[
k(m+n+1)<mn.
\]
For square $n\times n$:
\[
\text{ratio}\approx \frac{2k}{n}.
\]

\section{Nonstandard 5: Distinguish Two Quantum States}
Question type: ``Give a measurement basis where the probabilities differ.''
\begin{itemize}
\item First decide if the two expressions are actually the same state up to global phase. If yes, stop.
\item If not, choose a basis adapted to one of the states.
\item Best choices are often $\{\ket+,\ket-\}$ or $\{\ket{i},\ket{-i}\}$.
\end{itemize}
Example from exam:
\[
\ket{\psi_1}=\ket+,\qquad
\ket{\psi_2}=\frac{1}{\sqrt2}(-\ket0+i\ket1).
\]
Measure in $\{\ket+,\ket-\}$:
\[
P_{\psi_1}(\ket+)=1,
\qquad
P_{\psi_2}(\ket+)=\left|\frac{-1+i}{2}\right|^2=\frac12.
\]
Done.

\section{Nonstandard 6: Superposition Depends On Basis}
Question type: ``Is this state a superposition?''
Always ask: \textbf{with respect to which basis?}
\begin{itemize}
\item In standard basis, $\ket+$ and $\ket-$ are superpositions.
\item In Hadamard basis, $\ket+$ and $\ket-$ are basis vectors, hence not superpositions.
\item From the exam:
\[
\frac{1}{\sqrt2}(\ket+ + \ket-) = \ket0,
\]
so not a superposition in the standard basis.
\end{itemize}

\section{Usual Problem: 2x2 SVD By Hand}
\smallhead{Model problem.} Find an SVD of
\[
A=\begin{pmatrix}2&-1\\2&2\end{pmatrix}.
\]
\smallhead{Step 1.} Compute
\[
A^TA=\begin{pmatrix}8&2\\2&5\end{pmatrix}.
\]
\smallhead{Step 2.} Characteristic polynomial:
\[
\det(A^TA-\lambda I)
=
\begin{vmatrix}8-\lambda&2\\2&5-\lambda\end{vmatrix}
=\lambda^2-13\lambda+36.
\]
So
\[
\lambda_1=9,\qquad \lambda_2=4.
\]
Hence
\[
\sigma_1=3,\qquad \sigma_2=2.
\]
\smallhead{Step 3.} Right singular vectors from $A^TA$:
\[
v_1=\frac{1}{\sqrt5}\binom21,\qquad
v_2=\frac{1}{\sqrt5}\binom{1}{-2}.
\]
\smallhead{Step 4.} Left singular vectors:
\[
u_1=\frac{1}{3}Av_1=\frac{1}{\sqrt5}\binom12,\qquad
u_2=\frac{1}{2}Av_2=\frac{1}{\sqrt5}\binom{2}{-1}.
\]
\smallhead{Step 5.} Assemble
\[
U=\begin{pmatrix}
1/\sqrt5&2/\sqrt5\\
2/\sqrt5&-1/\sqrt5
\end{pmatrix},
\quad
\Sigma=\begin{pmatrix}3&0\\0&2\end{pmatrix},
\]
\[
V=\begin{pmatrix}
2/\sqrt5&1/\sqrt5\\
1/\sqrt5&-2/\sqrt5
\end{pmatrix},
\quad
A=U\Sigma V^T.
\]
\smallhead{Exam habit.} Always box:
\[
\sigma_i,\; v_i,\; u_i,\; U,\Sigma,V.
\]

\section{Usual Problem: Measurement Probability}
If
\[
\ket{\psi}=\frac{\sqrt3}{2}\ket0-\frac12\ket1
\]
and basis is $\{\ket0,\ket1\}$, read coefficients directly:
\[
P(\ket0)=\left|\frac{\sqrt3}{2}\right|^2=\frac34,
\qquad
P(\ket1)=\left|-\frac12\right|^2=\frac14.
\]
If basis is not standard, first rewrite state in that basis or compute inner products.

\section{Exam Checklist}
\begin{itemize}
\item Spectral question? First ask: is the matrix symmetric?
\item SVD question? Start with $A^TA$.
\item Rank/Col/Nul from SVD? Read from $\Sigma,U,V$; do not recompute.
\item Image compression? Quality uses squares of singular values.
\item Quantum same/different? Check global phase first.
\item Quantum probability? Basis vector inner products, then square magnitudes.
\item T/F proof? Write one defining identity, not a paragraph of fluff.
\end{itemize}

\section{Fast Memory Page}
\[
A=A^T \iff A=PDP^T.
\]
\[
A=U\Sigma V^T,\quad \sigma_i=\sqrt{\lambda_i(A^TA)}.
\]
\[
u_i=\frac{Av_i}{\sigma_i}.
\]
\[
\Rank(A)=\#\{\sigma_i>0\}.
\]
\[
A_k=\sum_{i=1}^k \sigma_i u_i v_i^T.
\]
\[
\text{quality}(k)=\frac{\sum_{i=1}^k \sigma_i^2}{\sum_i \sigma_i^2}.
\]
\[
\text{ratio}=\frac{k(m+n+1)}{mn}.
\]
\[
\ket+ = \frac{\ket0+\ket1}{\sqrt2},\quad
\ket- = \frac{\ket0-\ket1}{\sqrt2}.
\]
\[
\ket{i} = \frac{\ket0+i\ket1}{\sqrt2},\quad
\ket{-i} = \frac{\ket0-i\ket1}{\sqrt2}.
\]
\[
\ket{\psi}\sim e^{i\theta}\ket{\psi}.
\]
\[
P(\ket b)=|\braket{b}{\psi}|^2.
\]

\section{Usual Problem: Orthogonal Diagonalization}
Given symmetric
\[
A=\begin{pmatrix}a&b\\b&d\end{pmatrix},
\]
do:
\begin{itemize}
\item solve $\det(A-\lambda I)=0$,
\item find eigenvectors for each $\lambda$,
\item normalize them to $u_1,u_2$,
\item set
\[
P=[u_1\;u_2],\qquad D=\diag(\lambda_1,\lambda_2),
\]
\item conclude
\[
A=PDP^T.
\]
\end{itemize}
For repeated eigenvalue in symmetric case, eigenspace already has enough dimensions; just choose an orthonormal basis of it.

\section{Usual Problem: Read Rank/Col/Nul}
If
\[
\Sigma=\diag(\sigma_1,\sigma_2,0,0),
\]
then immediately:
\[
\Rank(A)=2,
\]
\[
\Col(A)=\Span\{u_1,u_2\},
\]
\[
\Nul(A)=\Span\{v_3,v_4\}.
\]
If a problem asks for a basis, write actual columns, not the symbols only.

\section{Usual Problem: Quality From Singular Values}
If singular values are $\sigma_1,\dots,\sigma_n$ and you keep first $k$, use
\[
\text{quality}(k)=\frac{\sum_{i=1}^k\sigma_i^2}{\sum_{i=1}^n\sigma_i^2}.
\]
For a \emph{percentage of singular values kept}, first compute $k$:
\[
k=\text{round}(pn)
\]
if the problem uses a fraction $p$ such as $0.75,0.50,0.25$.

\section{Usual Problem: Basis Expansion}
To express a state in a different basis, either use known identities or inner products.
Example:
\[
\ket0=\frac{1}{\sqrt2}\ket+ + \frac{1}{\sqrt2}\ket-.
\]
So in $\{\ket+,\ket-\}$ the coefficient vector is
\[
\binom{1/\sqrt2}{1/\sqrt2}.
\]
Likewise
\[
\ket1=\frac{1}{\sqrt2}\ket+ - \frac{1}{\sqrt2}\ket-.
\]

\section{Fast T/F One-Liners}
\begin{itemize}
\item Orthogonally diagonalizable $\Rightarrow$ symmetric.
\item Symmetric $\Rightarrow$ real eigenvalues and orthonormal eigenbasis.
\item Orthogonal does \emph{not} imply orthogonally diagonalizable.
\item $vv^T$ projector only if $v$ is unit.
\item Distinct eigenvalues of symmetric matrix $\Rightarrow$ orthogonal eigenvectors.
\item $\sigma_i^2$ are eigenvalues of $A^TA$.
\item $\sigma_i$ are never negative.
\item $\max_{\|x\|=1}\|Ax\|=\sigma_1$.
\item Same quantum state means same up to global phase.
\item Superposition is basis-dependent.
\end{itemize}

\end{multicols*}
\end{document}
